\textbf{COLUNI 2026 -- Questão 31}\par\medskip Leia o texto a seguir:

\textbf{É abril}

É abril na minha cidade.
É abril no mundo inteiro.
Sobe da terra tranquila um estímulo de vida e paz.
Um docel muito azul e muito alto cobre os reinos de Goiás.
Um sol de ouro novo vai virando e fugindo a longínquas partes do mundo.
Desaguaram em março as últimas chuvaradas do verão passado.

É festa alegre das colheitas.
Colhem-se as lavouras.
Quebra-se o milho maduro.
Bate-se o feijão,
já se cortou e se empilhou o arroz das roças.
As máquinas beneficiam o novo
e as panelas cozinham depressa o feijão novo e gostoso.
A abundância das lavouras são carreadas para depósitos e mercados.
Encostam-se nas máquinas os caminhões em carga completa.
Homens fortes, morenos, de dorso nu e reluzente,
descarregam e empilham a sacaria pesada.
Fecham-se quarteirões de ruas para a secagem de grãos
que secadores já não comportam.
Gira o capital, liquida-se nos bancos, paga-se no comércio. [...]
É um abril de bênçãos e aleluias
e cantam nas madrugadas todos os galos do mundo.
Os pássaros, os bichos se fartam nas sobras do que vai perdido pelas roças.
Respingam aqueles que não plantam nem colhem
e têm o direito às sobras dos que plantam e colhem. [...]
O leite transborda dos latões no rumo das cooperativas. [...]
E a vida se renova na força contagiante do trabalho.
Um sentido de fartura abençoa os reinos da minha cidade.\par
Disponível em: www.tudoepoema.com.br/cora-coralina-a-fala-de-aninha-e-abril/. Adaptado. Acesso em: 05 ago. 2025.

A partir do texto e dos conhecimentos geográficos, assinale a afirmativa que contém elementos que não se relacionam diretamente a fenômenos da natureza:\begin{enumerate}
\item As águas das chuvas e a fertilidade do solo.
\item Os ciclos da natureza e as espécies da fauna e flora.
\item A força de trabalho e as atividades econômicas.
\item As estações do ano e os períodos de plantio e colheita.\end{enumerate}